% ================================================================
% SECTION 4 — RESULTS
% Target: ~9 pages
% ================================================================

\section{Results}
\label{sec:results}

This section presents the empirical results in four stages.
Section~\ref{sec:res_eda} summarises the exploratory characteristics
of the panel data and the NATCOR series in particular.
Section~\ref{sec:res_1donor} demonstrates the consequences of an
inadequate donor pool using the single-corridor analysis from the
initial modelling iteration.
Section~\ref{sec:res_5donor} presents the main causal estimates
obtained with the full five-corridor donor pool, including pre-period
fit diagnostics and post-period gap series for each of the five
estimators.
Section~\ref{sec:res_robustness} presents the robustness battery:
in-time placebo tests, leave-one-out donor stability checks, and
confidence intervals derived from jackknife and block bootstrap
resampling.

% ---------------------------------------------------------------
\subsection{Exploratory Data Analysis}
\label{sec:res_eda}
% ---------------------------------------------------------------

The cleaned panel spans 120 months (January 2015 to December 2024),
six corridors, and 14 variables, yielding 720 observations with no
missing values after preprocessing.
Figure~\ref{fig:eda_all} plots monthly freight volume for all six
corridors over the full window.
Several features are immediately apparent.

\begin{figure}[H]
  \centering
  \includegraphics[width=\textwidth]{fig_eda_all_corridors}
  \caption{Monthly freight volume (million tonnes) for all six South
           African rail corridors, January 2015 -- December 2024.
           The grey band marks the COVID-19 disruption window
           (March--August 2020); the red dashed line marks the April
           2022 KZN flood event.
           Shading intensity reflects the severity of load-shedding
           stages recorded in each quarter.}
  \label{fig:eda_all}
\end{figure}

IRON\_ORE dominates throughput at 4.5--6\,MT per month, operating
largely independently of general freight dynamics owing to its
captive, export-commodity client base.
NATCOR and SOUTHCOR track each other at intermediate volume levels
(2--5\,MT range), consistent with their shared role as general freight
arteries serving the Gauteng industrial heartland.
CAPE, NE\_CORRIDOR, and NORTHCOR operate at lower volumes but exhibit
cyclical patterns broadly co-moving with NATCOR.
The COVID-19 disruption of March--August 2020 produced a sharp,
synchronised decline across all corridors, after which volumes
partially recovered before deteriorating again in 2021--2022 under
worsening load-shedding and cable theft pressure.
The April 2022 vertical marker coincides visually with a pronounced
inflection in NATCOR volume that is not matched in magnitude by any
other corridor.

Figure~\ref{fig:eda_natcor} provides a closer view of NATCOR's
behaviour, decomposing raw volume and the corridor's inferred market
share of total panel throughput.

\begin{figure}[H]
  \centering
  \includegraphics[width=\textwidth]{fig_eda_natcor_detail}
  \caption{NATCOR monthly freight volume (upper panel) and NATCOR
           market share as a fraction of total six-corridor throughput
           (lower panel), 2015--2024.
           The pre-flood mean and post-flood mean are shown as
           horizontal reference lines in the upper panel.}
  \label{fig:eda_natcor}
\end{figure}

In the pre-flood period (January 2015 to March 2022), NATCOR averaged
3.94\,MT per month with a standard deviation of 0.42\,MT.
Its market share was relatively stable at approximately 18--21\% of
total panel throughput.
Following the April 2022 shock, the post-period average falls to
3.63\,MT --- a raw decline of approximately 0.31\,MT per month --- and
market share contracts visibly.
Whether this observed decline is attributable to the flood itself, to
the concurrent deterioration in locomotive availability, or to broader
economic headwinds is precisely the question that the SCM estimators
are designed to answer.

The corridor-level correlations in the pre-period confirm the
suitability of the five donor corridors as controls.
Pearson correlations between NATCOR and each donor during 2015--2022
range from 0.61 (IRON\_ORE) to 0.83 (SOUTHCOR), and the first
principal component of the five-donor matrix explains 71\% of
cross-sectional variance.
These statistics support the assumption that a convex combination of
donor corridors can plausibly approximate the trajectory NATCOR would
have followed in the absence of the flood.

% ---------------------------------------------------------------
\subsection{Consequences of an Inadequate Donor Pool}
\label{sec:res_1donor}
% ---------------------------------------------------------------

Before presenting the main results, it is instructive to examine what
happens when the SCM framework is applied with only a single donor
corridor.
This scenario is not hypothetical: the initial modelling iteration
--- conducted before the full six-corridor panel was available ---
used only the CAPE corridor as the donor pool.
CAPE is a reasonable candidate on substantive grounds: it serves the
Western Cape export hinterland and shares some seasonal co-movement
with NATCOR.
However, a single donor unit cannot satisfy the core requirement of
synthetic control estimation, which is that the treated unit's
pre-period trajectory be reproducible as a convex combination of donor
units.
When only one donor is available, the ``combination'' is trivially
the donor itself, scaled by a single weight --- a degenerate case that
strips the method of its identification power.

Table~\ref{tab:1donor} reports the cumulative treatment effect
estimates produced by all five estimators when restricted to the
single-CAPE donor pool.

\begin{table}[H]
  \centering
  \caption{Cumulative post-flood treatment effect estimates (April
           2022 -- December 2024, 33 months) from the five SCM-family
           estimators using only the CAPE corridor as the donor pool.}
  \label{tab:1donor}
  \begin{tabularx}{\textwidth}{lXXX}
    \toprule
    Estimator & Pre-period RMSPE (MT) & Post-period RMSPE (MT)
              & Cumulative effect (MT) \\
    \midrule
    Base SCM  & 0.4917 & 0.8243 & $-11.104$ \\
    ASCM      & 0.4890 & 0.2703 & $+3.682$  \\
    EN-SCM    & 0.4921 & 0.6044 & $-1.458$  \\
    MC-SCM    & 0.0000 & 0.4791 & $-2.940$  \\
    SDID      & 0.4917 & 0.4461 & $-6.870$  \\
    \bottomrule
  \end{tabularx}
\end{table}

The estimates range from $-11.1$\,MT (Base SCM) to $+3.7$\,MT
(ASCM), a span of nearly 15\,MT.
The two most structurally different methods --- Base SCM and ASCM ---
disagree not only in magnitude but in sign.
Pre-period fit is poor across all methods (RMSPE $\approx 0.49$\,MT,
versus the 0.026\,MT achieved with five donors), confirming that CAPE
alone cannot reproduce NATCOR's pre-flood trajectory.

Figure~\ref{fig:donor_compare} juxtaposes the 1-donor and 5-donor
cumulative effect estimates to visualise the estimator divergence.

\begin{figure}[H]
  \centering
  \includegraphics[width=\textwidth]{fig_donor_comparison}
  \caption{Cumulative post-flood treatment effect (MT) estimated by
           the five SCM-family methods under the 1-donor (CAPE only)
           and 5-donor configurations.
           Error bars in the 5-donor panel show the range across
           leave-one-out subsets.}
  \label{fig:donor_compare}
\end{figure}

The contrast is stark.
With one donor the five methods are statistically and economically
irreconcilable; with five donors they collapse onto a consensus
band of just 0.15\,MT.
This finding provides an empirical demonstration of a theoretical
property of the SCM: estimator coherence is a consequence of adequate
donor pool coverage, not a feature of any individual estimation
algorithm.
Any one of the five methods, applied naively to the single-donor
panel, would have produced a result that could mislead policy analysis
by orders of magnitude.
The remainder of this section reports results from the five-donor
configuration exclusively.

% ---------------------------------------------------------------
\subsection{Main Causal Estimates: Five-Donor Configuration}
\label{sec:res_5donor}
% ---------------------------------------------------------------

\subsubsection{Pre-period Fit}

Table~\ref{tab:5donor} summarises the pre-period root mean squared
prediction error (RMSPE) and key post-period statistics for each of
the five estimators under the full five-corridor donor pool.

\begin{table}[H]
  \centering
  \caption{SCM-family estimator performance under the five-corridor
           donor pool (CAPE, IRON\_ORE, NE\_CORRIDOR, NORTHCOR,
           SOUTHCOR). Pre-period: January 2015 -- March 2022
           (87 months). Post-period: April 2022 -- December 2024
           (33 months).}
  \label{tab:5donor}
  \begin{tabularx}{\textwidth}{lXXXX}
    \toprule
    Estimator & Pre RMSPE (MT) & Post RMSPE (MT)
              & ATT (MT/month) & Cumulative (MT) \\
    \midrule
    Base SCM  & 0.0261 & 0.1704 & $-0.1582$ & $-5.2075$ \\
    ASCM      & 0.0260 & 0.1685 & $-0.1550$ & $-5.1109$ \\
    EN-SCM    & 0.0277 & 0.1703 & $-0.1582$ & $-5.2128$ \\
    MC-SCM    & 0.0000 & 0.1667 & $-0.1540$ & $-5.0834$ \\
    SDID      & 0.0261 & 0.1713 & $-0.1588$ & $-5.2401$ \\
    \midrule
    \textbf{Consensus}
              & \textbf{0.0212} & \textbf{0.1694}
              & $\mathbf{-0.1568}$ & $\mathbf{-5.1709}$ \\
    \bottomrule
  \end{tabularx}
\end{table}

Pre-period fit is excellent across all five methods.
RMSPE values of 0.026--0.028\,MT correspond to prediction errors of
less than 0.7\% of the NATCOR pre-period mean (3.94\,MT), confirming
that the synthetic control closely tracks NATCOR's historical
trajectory.
MC-SCM achieves a pre-period RMSPE of exactly zero by construction:
its nuclear norm regularisation problem is solved at a rank that
perfectly interpolates the pre-period observations.
The near-identical fit quality across the remaining four methods,
despite their structurally different estimation approaches, provides
initial reassurance that the synthetic donor combination is not
artefactually specific to a single algorithm.

\subsubsection{Synthetic Counterfactual and Gap Series}

Figure~\ref{fig:master_cf} plots the observed NATCOR volume alongside
the synthetic counterfactual produced by each estimator, together with
the implied gap (actual minus synthetic) for the post-flood period.

\begin{figure}[H]
  \centering
  \includegraphics[width=\textwidth]{fig_master_counterfactual}
  \caption{Upper panel: observed NATCOR monthly freight volume (black)
           versus synthetic counterfactuals from each of the five
           SCM-family estimators, 2015--2024.
           The red dashed vertical line marks the April 2022 flood.
           Lower panel: monthly gap series (actual minus synthetic) for
           each estimator in the post-flood period.
           The grey band spans the minimum-to-maximum envelope across
           all five methods; dashed lines show the individual method
           trajectories.}
  \label{fig:master_cf}
\end{figure}

Several observations follow from Figure~\ref{fig:master_cf}.
In the pre-period, all five synthetic counterfactuals track the
observed NATCOR series closely, confirming the pre-period fit
statistics.
The counterfactuals diverge slightly from one another owing to
differences in regularisation and weighting; however, these
differences are economically negligible (< 0.03\,MT).

In the post-period, the synthetic counterfactuals suggest NATCOR
should have been tracking at approximately 3.80--3.85\,MT per month
in the absence of the flood.
The observed series instead averages 3.63\,MT, producing a persistent
negative gap.
The gap series in the lower panel is uniformly negative across all 33
post-flood months for every estimator, with no month producing a
positive gap.
The largest negative gaps occur in the first six months following the
flood (May--October 2022), consistent with the initial infrastructure
disruption and emergency repair period.
Gaps then partially narrow but do not return to zero within the
observation window, suggesting that the flood effect has not fully
dissipated by December 2024.

\subsubsection{Estimator-Specific Weights}

The Base SCM solves for a closed-form constrained least-squares
solution over the pre-period.
The optimised donor weights assign the largest share to SOUTHCOR
($w = 0.463$), followed by NE\_CORRIDOR ($w = 0.241$), CAPE
($w = 0.221$), NORTHCOR ($w = 0.076$), and IRON\_ORE ($w = 0.000$).
The exclusion of IRON\_ORE from the convex combination is consistent
with its volume profile: IRON\_ORE operates at two to three times the
throughput of NATCOR and co-moves with export commodity prices rather
than with general freight demand.
The dominance of SOUTHCOR and NE\_CORRIDOR in the synthetic control
is economically intuitive, as both corridors serve the interior--coast
general freight market and exhibit similar seasonality patterns to
NATCOR.

ASCM begins with the Base SCM weights and applies a ridge regression
bias-correction (penalisation parameter $\alpha = 100$), which
produces a slight reweighting of the donor pool toward a more
diffuse allocation.
EN-SCM ($\alpha = 0.0005$, $\rho = 0.25$) selects a sparse weight
vector through $\ell_1$ penalisation while retaining a small
$\ell_2$ term, converging on a solution qualitatively similar to the
Base SCM but with somewhat stronger emphasis on SOUTHCOR.
MC-SCM decomposes the panel through singular value thresholding at
$\lambda = 0.2$ and recovers a rank-4 factor structure; its implicit
weighting is diffuse and not directly interpretable as corridor
weights.
SDID constructs both unit weights ($\hat{\omega}$) and time weights
($\hat{\lambda}$), upweighting pre-flood months in which donor
trajectories most closely resemble the post-period donor mean,
producing a quasi-matched pre-period window concentrated in 2019--2021.

% ---------------------------------------------------------------
\subsection{Robustness Analysis}
\label{sec:res_robustness}
% ---------------------------------------------------------------

\subsubsection{In-Time Placebo Tests}

To assess the statistical significance of the estimated flood effect,
in-time placebo tests were conducted for each estimator.
Five pseudo-intervention dates were used: January 2017, June 2018,
July 2019, April 2020, and January 2021.
For each placebo date $t^*$, the flood is artificially placed at $t^*$
and the post-period effect is estimated using the same methodology as
the true treatment.
An estimator is said to ``pass'' the placebo test if none of the five
placebo effects exceeds the true effect estimate in absolute magnitude.

\begin{table}[H]
  \centering
  \caption{In-time placebo test results for each SCM-family estimator.
           Columns report: the true cumulative treatment effect, the
           maximum absolute placebo effect across the five placebo
           dates, and the number of placebo exceedances (i.e.,
           placebo effects larger in magnitude than the true
           estimate).}
  \label{tab:placebo}
  \begin{tabularx}{\textwidth}{lXXX}
    \toprule
    Estimator & True effect (MT) & Max placebo (MT)
              & Exceedances / 5 \\
    \midrule
    Base SCM  & $-5.2075$ & $1.8340$ & 0 / 5 \\
    ASCM      & $-5.1109$ & $1.7921$ & 0 / 5 \\
    EN-SCM    & $-5.2128$ & $1.8651$ & 0 / 5 \\
    MC-SCM    & $-5.0834$ & $1.9203$ & 0 / 5 \\
    SDID      & $-5.2401$ & $1.8092$ & 0 / 5 \\
    \bottomrule
  \end{tabularx}
\end{table}

Every estimator achieves zero placebo exceedances (0/5).
The largest placebo effect in absolute value --- across all methods
and all five pseudo-dates --- is 1.92\,MT (MC-SCM at the 2018-01
placebo date), which is less than 38\% of the true MC-SCM estimate
of 5.08\,MT.
This result implies an exact Fisher $p$-value of $1/6 \approx 0.167$
under the permutation null (adding the true treatment to the five
placebos), which, while not meeting conventional $p < 0.05$ thresholds
in isolation, should be interpreted in the context of a donor pool of
only five units --- the maximum achievable permutation precision.
More importantly, the monotone ordering of the true effect over all
placebo effects is consistent across all five structurally distinct
estimators, providing convergent evidence that the flood produced a
genuinely exceptional decline in NATCOR volume.

Figure~\ref{fig:placebo} displays the placebo results graphically for
the ASCM, EN-SCM, and SDID estimators.

\begin{figure}[H]
  \centering
  \includegraphics[width=\textwidth]{fig_placebo_summary}
  \caption{In-time placebo test results for ASCM (left), EN-SCM
           (centre), and SDID (right).
           Each panel plots the cumulative gap series for the five
           pseudo-treatment dates (grey) and the true April 2022
           treatment (black).
           The true effect is the most negative trajectory in every
           panel.}
  \label{fig:placebo}
\end{figure}

\subsubsection{Leave-One-Out Donor Stability}

Leave-one-out (LOO) tests assess sensitivity to any individual donor
corridor by successively excluding each of the five donors from the
pool and re-estimating the treatment effect.
The LOO spread is defined as the range of cumulative effect estimates
across the five LOO subsets.

\begin{table}[H]
  \centering
  \caption{Leave-one-out donor stability results for the five
           SCM-family estimators.
           The LOO spread is the range (maximum minus minimum)
           of the cumulative effect across the five LOO subsets.
           A small spread indicates that no single donor corridor
           is driving the estimate.}
  \label{tab:loo}
  \begin{tabularx}{\textwidth}{lXXX}
    \toprule
    Estimator & Min LOO effect (MT) & Max LOO effect (MT)
              & LOO spread (MT) \\
    \midrule
    Base SCM  & $-5.235$ & $-5.155$ & 0.080 \\
    ASCM      & $-5.130$ & $-5.094$ & 0.036 \\
    EN-SCM    & $-5.261$ & $-5.171$ & 0.090 \\
    MC-SCM    & $-5.116$ & $-5.048$ & 0.068 \\
    SDID      & $-5.315$ & $-5.166$ & 0.149 \\
    \bottomrule
  \end{tabularx}
\end{table}

LOO spreads range from 0.036\,MT (ASCM) to 0.149\,MT (SDID), all
well below 3\% of the point estimates.
The most influential donor corridor across methods is SOUTHCOR: its
exclusion causes the largest downward revision to the estimated effect
(i.e., a smaller absolute loss), consistent with its role as the
highest-weighted corridor in the Base SCM solution.
Nonetheless, even excluding SOUTHCOR, every method continues to
estimate a cumulative loss in excess of $-5.0$\,MT, and the consensus
mean shifts by less than 0.10\,MT.
These results confirm that the causal estimate is not artefactually
dependent on any single donor corridor.

Figure~\ref{fig:loo} displays the LOO gap series for ASCM, EN-SCM,
and SDID.

\begin{figure}[H]
  \centering
  \includegraphics[width=\textwidth]{fig_loo_stability}
  \caption{Leave-one-out donor stability for ASCM (left), EN-SCM
           (centre), and SDID (right).
           Each panel overlays the five LOO gap series (coloured,
           one per excluded donor) against the full-pool estimate
           (black).
           Tight clustering of the LOO series around the full-pool
           estimate indicates low sensitivity to any individual donor.}
  \label{fig:loo}
\end{figure}

\subsubsection{Confidence Intervals}

Two interval estimation procedures are reported, one for each
estimator that supports it.

\textbf{SDID Jackknife Standard Error.}
The SDID monthly ATT estimate of $-0.1588$\,MT/month has a jackknife
standard error of $\widehat{\mathrm{se}}_{\mathrm{JK}} = 0.0024$\,MT,
corresponding to a $t$-statistic of $-65.6$ against the null of zero
effect.
The implied 95\% jackknife confidence interval for the cumulative
33-month effect is $[-5.393,\,-5.080]$\,MT.

\textbf{MC-SCM Block Bootstrap.}
A block bootstrap with block length 6 months and 500 replications
yields a 95\% interval for the MC-SCM cumulative estimate of
$[-5.210,\,-4.948]$\,MT.

Both intervals exclude zero by a wide margin.
Their respective widths (0.313\,MT for SDID, 0.262\,MT for MC-SCM)
are consistent with one another and small relative to the point
estimate, further confirming the precision of the causal identification.

\subsubsection{Consensus Estimate and Economic Magnitude}

Figure~\ref{fig:consensus} summarises the main result: the five-method
consensus estimate of the cumulative freight volume loss attributable
to the 2022 KZN floods.

\begin{figure}[H]
  \centering
  \includegraphics[width=\textwidth]{fig_consensus_effect}
  \caption{Consensus causal effect estimate. Each bar shows the
           cumulative 33-month treatment effect (MT) for one of the
           five SCM-family estimators under the five-corridor donor
           pool.
           Error bars for SDID and MC-SCM show 95\% confidence
           intervals (jackknife and block bootstrap, respectively).
           The dashed horizontal line marks the consensus mean of
           $-5.17$\,MT.}
  \label{fig:consensus}
\end{figure}

Across the five estimators, cumulative effect estimates range from
$-5.08$\,MT (MC-SCM) to $-5.24$\,MT (EN-SCM), with a consensus mean
of $-5.17$\,MT and a total spread of 0.15\,MT.
The average monthly effect is $-0.155$ to $-0.159$\,MT/month,
sustained uniformly across all 33 post-flood months.

To contextualise the economic magnitude, consider that NATCOR's
pre-flood mean was 3.94\,MT/month.
A sustained loss of $\approx 0.157$\,MT/month represents a 4.0\%
reduction in corridor throughput.
At an estimated rail freight revenue rate of approximately
ZAR\,250 per tonne (consistent with Transnet's published tariff
schedules for general freight), the cumulative revenue loss attributable
directly to the flood effect is approximately
ZAR\,1.28 billion over 33 months.
This figure does not include secondary economic costs --- road damage
from diverted heavy freight, shipper inventory accumulation costs,
or the modal shift to road that may be partially permanent --- which
would substantially increase the total welfare cost of the disruption.

The fact that the gap series remains uniformly negative through
December 2024 --- 33 months after the flood --- implies that NATCOR
had not fully recovered to its counterfactual trajectory by the end
of the observation window.
This persistence is consistent with a structural demand-side effect,
in which shippers who transferred freight to road transport during the
disruption period did not fully return to rail after infrastructure
repairs were completed.
Section~\ref{sec:discussion} returns to this interpretation and its
policy implications.
